% =========================================================
% Scratch: No 3-Element Boolean Ring — Example Note
% Chapter: Boolean Rings and Boolean Algebras — Volume I
% Source: Halmos & Givant, Introduction to Boolean Algebras, §1
%
% STATUS: scratch file — move into chapter structure when
%         volume-i/boolean-rings/ directory is created.
%
% Depends on:
%   \ref{lem:self-additive-inverse}  (p + p = 0)
%   colors: defbox/defborder, propbox/propborder (common/colors.tex)
%   environments: example*, remark* (common/environments.tex)
%   packages: booktabs (for \toprule etc.)
% =========================================================

% ---------------------------------------------------------
% Remark: No Boolean Ring of Order 3 Exists
% ---------------------------------------------------------

\begin{remark*}[No Boolean ring of order 3]
A Boolean ring of order~3 does not exist.
The additive group $(R,+)$ of any Boolean ring satisfies
$p + p = 0$ for all $p$, by Lemma~\ref{lem:self-additive-inverse}.
Every element therefore has additive order dividing~2, so
$(R,+)$ is an elementary abelian $2$-group.
Elementary abelian $2$-groups have order $2^n$ for some $n \ge 0$.
Since $3$ is not a power of~$2$, no Boolean ring of order~$3$
can exist.
\end{remark*}

% ---------------------------------------------------------
% Example: The Four-Element Boolean Ring 2^2
% ---------------------------------------------------------

\begin{example*}[The four-element Boolean ring $\mathbf{2}^2$]
\label{ex:boolean-ring-four-elements}
The smallest Boolean ring with more than two elements is
$\mathbf{2}^2 = \{(0,0),\,(0,1),\,(1,0),\,(1,1)\}$,
the set of ordered pairs from $\mathbf{2} = \{0,1\}$.
Addition and multiplication are defined coordinatewise:
\[
  (p_0,p_1) + (q_0,q_1) \;=\; (p_0+q_0,\;p_1+q_1),
  \qquad
  (p_0,p_1) \cdot (q_0,q_1) \;=\; (p_0 \cdot q_0,\;p_1 \cdot q_1),
\]
where each coordinate is computed in $\mathbf{2}$.
The zero element is $(0,0)$ and the unit is $(1,1)$.

\medskip
\noindent\textbf{Addition table.}
\begin{center}
\renewcommand{\arraystretch}{1.25}
\begin{tabular}{c|cccc}
$+$     & $(0,0)$ & $(0,1)$ & $(1,0)$ & $(1,1)$ \\ \hline
$(0,0)$ & $(0,0)$ & $(0,1)$ & $(1,0)$ & $(1,1)$ \\
$(0,1)$ & $(0,1)$ & $(0,0)$ & $(1,1)$ & $(1,0)$ \\
$(1,0)$ & $(1,0)$ & $(1,1)$ & $(0,0)$ & $(0,1)$ \\
$(1,1)$ & $(1,1)$ & $(1,0)$ & $(0,1)$ & $(0,0)$
\end{tabular}
\end{center}

\medskip
\noindent\textbf{Multiplication table.}
\begin{center}
\renewcommand{\arraystretch}{1.25}
\begin{tabular}{c|cccc}
$\cdot$ & $(0,0)$ & $(0,1)$ & $(1,0)$ & $(1,1)$ \\ \hline
$(0,0)$ & $(0,0)$ & $(0,0)$ & $(0,0)$ & $(0,0)$ \\
$(0,1)$ & $(0,0)$ & $(0,1)$ & $(0,0)$ & $(0,1)$ \\
$(1,0)$ & $(0,0)$ & $(0,0)$ & $(1,0)$ & $(1,0)$ \\
$(1,1)$ & $(0,0)$ & $(0,1)$ & $(1,0)$ & $(1,1)$
\end{tabular}
\end{center}

\medskip
\noindent\textbf{Verification of idempotence.}
Every element squares to itself:
\begin{center}
\renewcommand{\arraystretch}{1.2}
\begin{tabular}{cc}
\toprule
$p$     & $p \cdot p$ \\ \midrule
$(0,0)$ & $(0,0)$     \\
$(0,1)$ & $(0,1)$     \\
$(1,0)$ & $(1,0)$     \\
$(1,1)$ & $(1,1)$     \\ \bottomrule
\end{tabular}
\end{center}

\medskip
\noindent\textbf{Verification of characteristic~$2$.}
Every element is its own additive inverse:
\begin{center}
\renewcommand{\arraystretch}{1.2}
\begin{tabular}{cc}
\toprule
$p$     & $p + p$ \\ \midrule
$(0,0)$ & $(0,0)$ \\
$(0,1)$ & $(0,0)$ \\
$(1,0)$ & $(0,0)$ \\
$(1,1)$ & $(0,0)$ \\ \bottomrule
\end{tabular}
\end{center}
\end{example*}

% ---------------------------------------------------------
% Remark: The Pattern — Boolean Rings Have Order 2^n
% ---------------------------------------------------------

\begin{remark*}[Boolean rings have order $2^n$]
Every finite Boolean ring is isomorphic to $\mathbf{2}^n$ for
some $n \ge 0$, and therefore has exactly $2^n$ elements.
This follows from the representation theorem
(Halmos--Givant, Chapter~22): every Boolean algebra is isomorphic
to a field of sets, and every finite field of sets is isomorphic
to $\mathcal{P}(X)$ for some finite set $X$ with $|X| = n$,
giving $|\mathcal{P}(X)| = 2^n$.
The permissible finite sizes are $1, 2, 4, 8, 16, \ldots$
and no other cardinality is achievable.
\end{remark*}
